\subsection*{Compose Two Bounded Variables for Tiling with \texttt{\#}}

The example combines \textbf{two bounded variables} to select tiles. The \textbf{\texttt{\#} operator} composes their coordinates over a Cartesian product. Operand order matters: the operation is \textbf{not commutative}.

\noindent\textbf{The \texttt{\#} Operator:}\par

\begin{itemize}[leftmargin=1.6em]
\item The code defines the tiling parameters with four bounded variables: \texttt{q}, \texttt{n}, \texttt{h}, and \texttt{w}, whose upper bounds are \texttt{6}, \texttt{N/6}, \texttt{H/16}, and \texttt{W/512}.
\item Suppose \texttt{N} is 18, \texttt{q} represents physical thread index. You want each thread to process data that is contiguous in the first dimension. That is, the first thread handles input[0\textasciitilde{}2][xxx], the second thread handles input[3\textasciitilde{}5][xxx], and so on.
\item \texttt{q \# n} will result in a new bounded variable implicitly, whose upper bound is \texttt{N}: bound of \texttt{q} multiply bound of \texttt{n}.
\item For every data move, the stride of each dimension is 1, 3, 16, 512. In the last dimension, the index is \texttt{w}. In the first dimension, the index is \texttt{q * (bound of n) + n}.
\end{itemize}

\noindent\textbf{Why this is Important:}\par

\begin{itemize}[leftmargin=1.6em]
\item The combination of \texttt{q} and \texttt{n} ensures that physical and virtual parallelism can work together efficiently.
\item This helps optimize \textbf{memory access patterns}, especially when working with large matrices, where efficient memory use is crucial.
\item Croqtile's approach here is more aligned with a \textbf{physical storage hierarchy} (local, shared, and global storage). Unlike higher-level loop scheduling techniques that abstract away the actual data movement, Croqtile explicitly links \textbf{virtual loops} to \textbf{physical DMA operations}. This results in more efficient data movement and computation.
\end{itemize}
